\section{Large exact messages from fixed local data}\label{sec:message-limits}
Section~\ref{sec:hardness} shows that exact optimization can be hard.
The size of a message is a different obstacle, and it appears even when
optimization is easy. A message is the pointwise minimum of its
branches, and an exact dictionary must keep every branch that is the
unique minimum at some boundary value.

Section~\ref{sec:limits-chain} builds the instances of
Theorem~\ref{thm:fixed-data} from identical stages, and
Section~\ref{sec:limits-all} proves that one of their messages needs
$2^n$ branches in every exact dictionary, while $w$, $C$, $H$,
$\mu^{-1}$, and $R$ stay bounded and the input length is polynomial in
$n$. Hence, without perturbation, the dictionary size has no bound
polynomial in the input length and these quantities.
Section~\ref{sec:limits-pruning} treats approximation within a uniform
error $\varepsilon$ (Theorem~\ref{thm:pruning}).

\subsection{A chain that writes the support into its final state}
\label{sec:limits-chain}
Fix a rational $0<\theta\le1/10$ and define, for $t\in[0,1]$,
\begin{equation}\label{eq:constant-message}
 V_n(t)=\min\left\{
 \sum_{i=1}^n(x_i^2-2x_i+z_i)
 +\sum_{i=1}^n(s_i-\theta s_{i-1}-\theta x_i)^2:
 \begin{array}{l}
 x_i(1-z_i)=0,\quad z_i\in\{0,1\},\\ s_0=0,\quad s_n=t
 \end{array}\right\}.
\end{equation}
The controls $x_i$ have indicators $z_i$, and the states $s_i$ are free.
Each stage penalizes deviation from the dynamics
$s_i=\theta s_{i-1}+\theta x_i$, and every stage has the same
coefficients. On indicator-feasible points,
$x_i^2-2x_i+z_i=(x_i-z_i)^2$, so an inactive control is zero and an
active control costs nothing exactly when it equals one. The terminal
equality $s_n=t$ makes $V_n$ a conditional value function of the
boundary value $t$; it is not a constraint of the global problem. We
identify a support with its indicator vector $z\in\{0,1\}^n$. As in
Section~\ref{sec:prelim-messages}, the branch of $z$ is the minimum over
the continuous variables with the support fixed to $z$. Unlike the
message branches \eqref{eq:branch}, these branches keep the terms that
involve only $t$, which makes them convex rather than concave in $t$.
Up to a shift of $t$, the message constructed in the proof of
Theorem~\ref{thm:fixed-data} differs from $V_n$ by one fixed
polynomial, which does not change the number of formulas needed.

For a fixed support $z$, the choice $x=z$ with exact dynamics costs
nothing and drives the state to the \emph{center}
$p_z=\sum_{i=1}^n\theta^{n-i+1}z_i$. Each stage multiplies the state by
$\theta$, so the control of stage $i$ reaches the end of the chain with
weight $\theta^{n-i+1}$. The center is therefore the number with digits
$z_n,z_{n-1},\ldots,z_1$ in base $1/\theta$. For $\theta=1/10$ these
are ordinary decimal digits: with $n=3$, the support $z=(1,0,1)$ has
center $0.101$, and $z=(0,1,1)$ has center $0.110$. Reaching another
terminal value $t$ costs $(t-p_z)^2$ divided by a factor $D_z\in[1,2)$
(Lemma~\ref{lem:residual-message}). Because $\theta$ is small, distinct
digit strings give centers at least $\theta^n$ apart, so each support
is the unique best one near its own center.

This is how fixed local data produce exponentially many branches. No
coefficient is large, and none changes from stage to stage. What grows
with $n$ is the number of digits of $t$ that matter: to tell all $2^n$
supports apart, one must know $t$ to within about $\theta^n$, and a
single continuous boundary value can carry that many digits.

\begin{lemma}\label{lem:residual-message}
Let $W_n=\sum_{j=0}^{n-1}\theta^{2j}$. The branch of support $z$ in
\eqref{eq:constant-message} is
\begin{equation}\label{eq:constant-branches}
 q_z(t)=\frac{(t-p_z)^2}{D_z},\qquad
 p_z=\sum_{i=1}^n\theta^{n-i+1}z_i,\qquad
 D_z=W_n+\sum_{i=1}^n\theta^{2(n-i+1)}z_i.
\end{equation}
Thus $V_n=\min_z q_z$.
\end{lemma}
\begin{proof}
Fix $z$ and put $e_i=x_i-z_i$ and
$r_i=s_i-\theta s_{i-1}-\theta x_i$. Inactive controls force $e_i=0$,
and the objective is $\sum_i(e_i^2+r_i^2)$. Unrolling the dynamics from
$s_0=0$ to $s_n=t$ leaves the single linear constraint
\[
 t-p_z=\sum_i\theta^{n-i+1}e_i+\sum_i\theta^{n-i}r_i.
\]
Conversely, any residuals that satisfy it, with $e_i=0$ for inactive
$i$, determine feasible controls and states. The free residuals are
$e_i$ for active $i$ and every $r_i$, and their coefficients in this
constraint have squared norm $D_z$. Cauchy--Schwarz therefore gives the
minimum $(t-p_z)^2/D_z$, attained at
$e_i=(t-p_z)\theta^{n-i+1}z_i/D_z$ and
$r_i=(t-p_z)\theta^{n-i}/D_z$.
\end{proof}

\subsection{All branches are needed}
\label{sec:limits-all}
\begin{theorem}\label{thm:fixed-data}
Every finite exact representation of $V_n$ on $[0,1]$ as a minimum of
polynomials, or as a piecewise polynomial function, uses at least $2^n$
distinct polynomial formulas. The same holds for a scalar conditional
message, on any box $[-R,R]$ with $R\ge5$, of an unperturbed instance of
\eqref{eq:problem} with unit penalties on every variable, coefficients
from a fixed finite alphabet of rationals depending only on $\theta$,
and the graph and spectral bounds of Theorem~\ref{thm:hardness} (width
at most two for all $n$, and exactly two when $n\ge2$).
\end{theorem}
\begin{proof}
Each $q_z$ is the unique minimum on an interval around $p_z$, so every
finite representation must contain $q_z$ itself. For the second claim,
we give the free states unit-penalty indicators that no minimizer
switches off.

Put $\kappa_\theta=(1+\theta^2)/(1-\theta^2)<2$. Every $D_z$ lies in
$[1,\kappa_\theta]$, and all centers lie in $[0,\theta/(1-\theta)]$.
Distinct centers are at least $\theta^n$ apart. Indeed, the bit $z_i$
carries the power $\theta^{n-i+1}$. Let $k$ be the smallest exponent at
which two supports differ, that is, the leading digit in which their
centers differ. If $k<n$, their centers differ by at least
$\theta^k(1-2\theta)/(1-\theta)\ge\theta^n$, using $\theta\le1/10$;
if $k=n$, they differ by exactly $\theta^n$. On $|t-p_z|\le\theta^n/3$,
intersected with $[0,1]$, we have
\[
 q_z(t)\le\theta^{2n}/9,
 \qquad q_{z'}(t)\ge4\theta^{2n}/(9\kappa_\theta)>q_z(t)
 \quad(z'\ne z).
\]
So $q_z$ is the unique minimum, and equals $V_n$, on a nonempty open
interval. If finitely many polynomials represent $V_n$ there, one of
them agrees with $q_z$ at infinitely many points and hence identically,
because a nonzero polynomial has only finitely many roots. This
argument applies both to a minimum and to a piecewise representation,
and to polynomials of any degree. Distinct centers give distinct
polynomials.

It remains to realize $V_n$ as a message of \eqref{eq:problem}. Before
$s_n$ is fixed, the quadratic matrix in \eqref{eq:constant-message}, in
the order $(x_1,s_1,\ldots,x_n,s_n)$, has the edges
\eqref{eq:hard-edges}, diagonal entries $1+\theta^2$ for controls and
nonterminal states, and final diagonal entry one. The graph is that of
Theorem~\ref{thm:hardness}, and the same Gershgorin row bounds prove
\eqref{eq:hard-spectrum}; the final row has off-diagonal sum $2\theta$.

Problem \eqref{eq:problem} attaches an indicator to every variable,
while the states in \eqref{eq:constant-message} are free. We therefore
shift the states by four and give each one a unit-penalty indicator.
The shift moves the reference trajectories away from zero, so switching
a state indicator off, which forces that state to zero, leaves a large
residual. Let $y_0=4$, attach an indicator $\zeta_i$ to $y_i$, and
replace the dynamic terms by
\begin{equation}\label{eq:message-lift}
 \sum_i[y_i-\theta y_{i-1}-\theta x_i-4(1-\theta)]^2
       +\sum_i\zeta_i,
 \qquad y_n=4+t.
\end{equation}
With $y=4\mathbf1+s$ and every $\zeta_i=1$, the lifted objective is the
objective of \eqref{eq:constant-message} plus $n$. The separator is
$S=\{y_n\}$; the lower bound already holds on
$y_n\in[4,5]\subseteq[-R,R]$, which corresponds to $t\in[0,1]$.

No minimizer switches a state indicator off. To see this, fix a control
support $z$ with reference states $s_0^z=0$ and
$s_i^z=\theta s_{i-1}^z+\theta z_i\ge0$. With $e=x-z$ and
$d=y-4\mathbf1-s^z$, the control and dynamic squares are the left side
of \eqref{eq:activation-quadratic} with $b_i=\theta$, so their sum is
at least $(1-3\theta)(\|e\|^2+\|d\|^2)$. An indicator $\zeta_i=0$ forces
$y_i=0$, hence $d_i\le-4$. If $k$ state indicators are off, the
objective is at least $n-k+16(1-3\theta)k\ge n+10.2k$. The choice with
all state indicators on, $x=z=0$, $y_i=4$ for $i<n$, and $y_n=4+t$
costs $n+t^2\le n+1$. Hence every minimizer has all state indicators
on, and the conditional value is $n+V_n(t)$.

The message convention \eqref{eq:message} omits the terminal penalty,
which leaves $n-1+V_n(t)$. It also omits the terms involving only
$y_n$, and writing the lifted objective in the form \eqref{eq:problem}
drops an additive constant, which is at most $16n$. So the message
differs from $n-1+V_n$ by a fixed polynomial. Subtracting one fixed
polynomial from every formula does not change the number of
indispensable formulas. The omitted terms include $y_n^2$ with
coefficient one, while each $q_z$, as a function of $y_n$, has leading
coefficient $1/D_z\le1$; this is why the convex functions $q_z$ become
concave branches of the message. The expanded linear coefficients have
magnitude at most eight. The quadratic and linear coefficients and the
unit penalties belong to a finite alphabet depending only on $\theta$.
\end{proof}

The theorem counts explicit polynomial formulas. It is not a lower
bound for implicit representations, such as the optimization problem
\eqref{eq:constant-message} that defines $V_n$, nor for optimization,
because the global problem is trivial. Without the terminal condition,
every control support has the zero-residual trajectory $x=z$, so the
optimal value of \eqref{eq:constant-message} without $s_n=t$ is zero,
and that of the lifted problem is $n$. The example is thus an extreme
case of near ties between supports. Margin conditions, such as that of
\cite[Definition~5]{BhathenaStructured}, are designed to limit near ties
of this kind. The perturbation in Theorem~\ref{thm:main} shifts the
branch of each support $A$ by $\sum_{i\in A}\xi_i$, and these random
shifts limit near ties in expectation (Lemma~\ref{lem:count}); the
graph and the conditioning alone do not.

\subsection{Matching bounds for pruning branches}
\label{sec:limits-pruning}
If a uniform error $\varepsilon$ is allowed, a natural approach keeps
only some branches $q_z$ and discards the rest. Let $N_n(\varepsilon)$
be the smallest number of branches $q_z$ whose minimum $\widehat V$
satisfies $0\le\widehat V-V_n\le\varepsilon$ on $[0,1]$, and set
$N_\theta(\varepsilon)=\sup_{n\ge1}N_n(\varepsilon)$. This supremum is
finite, and Theorem~\ref{thm:pruning} gives its order of growth as
$\varepsilon\downarrow0$.

\begin{theorem}\label{thm:pruning}
As $\varepsilon\downarrow0$,
\[
 N_\theta(\varepsilon)=\Theta_\theta(\varepsilon^{-\alpha_\theta}),
 \qquad \alpha_\theta=\frac{\log2}{2\log(1/\theta)}.
\]
The lower bound concerns only minima of branches $q_z$.
\end{theorem}

\emph{Proof idea.} Keeping the $m$ most significant digits of the
centers (the bits $z_{n-m+1},\ldots,z_n$) locates every center to within
about $\theta^m$, which gives error of order $\theta^{2m}$ with about
$2^m$ branches. Conversely, $V_n$ vanishes at $2^m$ centers spaced at
least $\theta^m$ apart, and once $\varepsilon$ is small compared with
$\theta^{2m}$, each of these zeros needs its own branch. Each additional
digit thus doubles the number of branches and multiplies the error by
about $\theta^2$, which gives the exponent $\alpha_\theta$; for
$\theta=1/10$, $\alpha_\theta\approx0.15$.

\begin{proof}
For the upper bound, keep the last $m$ bits $z_{n-m+1},\ldots,z_n$ of a
support, where $0\le m\le n$; they carry the powers
$\theta,\ldots,\theta^m$. Let
\[
 \delta=\sum_{j=m+1}^n\theta^j\le\frac{\theta^{m+1}}{1-\theta},
 \qquad
 \beta=\sum_{j=m+1}^n\theta^{2j}\le\frac{\theta^{2m+2}}{1-\theta^2}.
\]
For each choice of the kept bits, retain the two supports whose omitted
bits are all zero or all one. These are at most $2^{m+1}$ of the
branches $q_z$, each with its own denominator. For an optimal support
$z$ at $t$, the two retained centers $p_0,p_1$ with the same kept bits
bracket $p_z$ and are $\delta$ apart. The one closer to $t$ has
squared distance at most $(t-p_z)^2+\delta^2/4$ from $t$: outside
$[p_0,p_1]$ it is no farther than $p_z$, and inside the interval it is
within $\delta/2$. Its denominator differs from $D_z$ by at most
$\beta$; all denominators are at least one and $(t-p_z)^2\le1$.
Consequently
\begin{equation}\label{eq:pruning-upper}
 0\le\widehat V-V_n\le\delta^2/4+\beta
 \le A_\theta\theta^{2m+2},\qquad
 A_\theta=\frac1{4(1-\theta)^2}+\frac1{1-\theta^2}.
\end{equation}
Taking $m=\lceil\log(A_\theta\theta^2/\varepsilon)/
(2\log(1/\theta))\rceil$, or retaining every support if $n\le m$,
proves the upper bound $4(A_\theta\theta^2/\varepsilon)^{\alpha_\theta}$
for sufficiently small $\varepsilon$.

For the lower bound take $n\ge m\ge1$. The $2^m$ supports whose first
$n-m$ bits are zero have centers separated by at least $\theta^m$, by
the separation argument in the proof of Theorem~\ref{thm:fixed-data},
and $V_n$ vanishes at each of these centers. Since
$\widehat V\le V_n+\varepsilon$, some retained branch is at most
$\varepsilon$ at each such center, so its own center lies within
$\sqrt{\kappa_\theta\varepsilon}$. If
$\varepsilon<\theta^{2m}/(4\kappa_\theta)$, a retained branch can serve
at most one of the $2^m$ centers. Choose
$m=\lfloor\log(1/(8\kappa_\theta\varepsilon))/(2\log(1/\theta))\rfloor$.
Then $N_\theta(\varepsilon)\ge
\tfrac12(8\kappa_\theta\varepsilon)^{-\alpha_\theta}$.
\end{proof}

Section~\ref{sec:related} compares Theorems~\ref{thm:fixed-data}
and~\ref{thm:pruning} with known value functions that have
exponentially many pieces \cite{Murty,KuricTree} and with known pruning
rates \cite{GaubertPruning}.
